\section{Process Supervision: verificación paso a paso del razonamiento}

\subsection{Diferencia entre Outcome y Process}

Supongamos una cadena de razonamiento con pasos $s_{1:T}$. Outcome supervision solo asigna una etiqueta $R\in\{0,1\}$ a la trayectoria completa; process supervision proporciona $r_t\in\{0,1\}$ o una puntuación continua para cada paso.

\videofigure{figures/outcome-supervision.jpg}{Outcome supervision solo entrega una etiqueta al final de todo el razonamiento.}{00:22:20--00:24:00}

\videofigure{figures/process-supervision.jpg}{Process supervision localiza con precisión el paso erróneo y mejora la asignación de crédito.}{00:24:00--00:26:02}

\begin{importantbox}{Dos valores de process supervision}
Primero, permite podar en cuanto aparece un error, ahorrando tokens posteriores; segundo, reduce los falsos positivos de "un proceso erróneo que por casualidad llega a la respuesta correcta". Para el entrenamiento, el reward a nivel de paso también ofrece una señal de aprendizaje más densa que un único reward final.
\end{importantbox}

\subsection{PRM800K: etiquetas de pasos hechas por humanos}

Let's Verify Step by Step recopila cerca de 800,000 etiquetas de pasos hechas por humanos. Los anotadores leen la solución generada por el modelo y, después de cada paso, juzgan si el proceso hasta ese momento sigue siendo correcto. Para mejorar la eficiencia de los datos, el método usa active learning, dando prioridad a los humanos en las muestras más discriminativas o con mayor probabilidad de revelar debilidades del modelo.

\videofigure{figures/prm-labeling.jpg}{PRM800K realiza anotación humana a nivel de paso sobre razonamientos matemáticos generados por el modelo.}{00:26:02--00:29:08}

Después del entrenamiento, el PRM produce la probabilidad de corrección $q_t$ de cada paso. La trayectoria completa puede agregarse mediante el paso más débil, el producto o la suma de logaritmos, por ejemplo:

$$
S(y)=\sum_{t=1}^{T}\log q_t.
$$

\begin{itemize}
    \item $q_t=P(s_t\text{ correct}\mid x,s_{1:t})$;
    \item $S(y)$: puntuación de proceso global de la trayectoria;
    \item la suma de logaritmos penaliza cualquier paso erróneo de baja probabilidad, y a la vez evita el desbordamiento por defecto de multiplicar directamente.
\end{itemize}

\subsection{Ganancias del PRM frente al ORM}

\videofigure{figures/prm-vs-orm.jpg}{El PRM supera al ORM y al majority voting en la selección de candidatos.}{00:29:08--00:30:44}

Las ganancias del PRM son especialmente notorias en cadenas de razonamiento largas, porque cuanto más larga es la trayectoria, mayor es la probabilidad de que un error local quede oculto por la respuesta final. También permite hacer beam search: se evalúa cada segmento generado y solo se expanden los prefijos con puntuación alta.

\subsection{Active Learning y anotación a gran escala}

El muestreo aleatorio desperdicia gran cantidad de esfuerzo humano en pasos claramente correctos o claramente incorrectos. Active learning selecciona las muestras en las que el verifier está actualmente más incierto, donde hay mayor discrepancia entre modelos o que están más cerca de la frontera de decisión.

\videofigure{figures/active-learning.jpg}{Active learning mejora el rendimiento por unidad de esfuerzo humano en la anotación de pasos.}{00:30:44--00:33:02}

\subsection{Generalización fuera de distribución}

\videofigure{figures/prm-ood.jpg}{El PRM entrenado conserva una ventaja relativa incluso en datos matemáticos fuera de distribución.}{00:33:02--00:35:20}

Esta generalización indica que el PRM aprendió cierta razonabilidad general de los pasos matemáticos, pero no debe interpretarse como que domina la demostración rigurosa. Los supuestos implícitos, los saltos de pasos y la ambigüedad simbólica del razonamiento en lenguaje natural aún pueden engañar al verifier.

\begin{warningbox}{El process reward también puede producir nuevo reward hacking}
El modelo puede aprender a escribir trayectorias extensas en las que "cada paso parece de libro de texto" para complacer al PRM, en lugar de resolver el problema de forma más eficaz. Si el verifier favorece cierto formato, longitud o fraseo común, el RL amplificará esas preferencias. La supervisión a nivel de paso aumenta la resolución, pero no elimina el riesgo del objetivo indirecto.
\end{warningbox}

\subsection{Resumen de la sección}

La process supervision hecha por humanos proporciona una señal de verificación precisa, densa y utilizable en búsqueda, al costo de una anotación experta costosa. La siguiente pregunta natural es: ¿se puede lograr que el modelo y el entorno generen etiquetas de paso de forma automática?

